% =====================================================================
%  CHECKPOINT 03 -- state after block 3 (maths)
%  Copy this over latex-workshop.tex if you need to catch up.
%
%  House style in these files: one sentence per line, \label on the
%  same line as its \section, four spaces per indent.
% =====================================================================

\documentclass[a4paper,11pt]{article}

% ---------------------------------------------------------------------
%  Preamble
% ---------------------------------------------------------------------
\usepackage[T1]{fontenc}        % proper 8-bit output encoding
\usepackage[english]{babel}     % hyphenation and language conventions
\usepackage{newtxtext}          % Times for the text
\usepackage{parskip}            % blank line between paragraphs, no indent
\usepackage{amsmath}            % align, cases, \text, decent spacing
\usepackage{newtxmath}          % ...and Times for the maths to match

% TODO 4.2  Add \usepackage{booktabs} and \usepackage{caption} here
%           once you get to the table.

% TODO 4.3  Add \usepackage{graphicx} here once you get to the figure.

% TODO 4.5  Add the bibliography machinery here:
%
%             \usepackage[backend=biber,style=numeric]{biblatex}
%             \addbibresource{references.bib}
%
%           On Overleaf that is all you need. Locally, use latexmk so
%           that biber is run for you.

\title{How Fast Does Your Coffee Go Cold?}
\author{Your Name}
\date{\today}

\begin{document}

\maketitle
\tableofcontents

% =====================================================================
\section{Introduction}\label{sec:introduction}

A freshly poured coffee is too hot to drink, and a forgotten one is too cold.
Somewhere in between there is a window of a few minutes in which the drink is actually pleasant,
and this report sets out to find it.
The underlying physics is \emph{Newton's law of cooling},
which says that a body loses heat in proportion to how much hotter it is than its surroundings.
Three \textbf{everyday vessels} are compared,
referred to throughout by the short labels \texttt{mug}, \texttt{glass} and \texttt{flask}.

The question is worth asking precisely because the model is so thin.
It has exactly one free parameter per vessel,
so a single afternoon is enough to test it properly.

This report sets out to:

\begin{itemize}
    \item state the cooling model and what its one parameter means,
    \item compare the cooling curves of the three vessels, and
    \item work out how long each of them keeps a drink in the pleasant range.
\end{itemize}

% =====================================================================
\section{Theory}\label{sec:theory}

The model is described by three quantities:

\begin{description}
    \item[$T_0$] temperature of the drink at the start
    \item[$T_\mathrm{env}$] temperature of the room
    \item[$\tau$] time constant of the vessel: the time it takes to close
        roughly two thirds of the gap to room temperature
\end{description}

Newton's law of cooling states that the temperature difference to the surroundings decays
exponentially, so that

\begin{equation}
    T(t) = T_\mathrm{env} + \left( T_0 - T_\mathrm{env} \right) \mathrm{e}^{-t/\tau} .
    \label{eq:cooling}
\end{equation}

Rearranging for the time gives the waiting time until the drink has reached a chosen
temperature $T$,

% \notag suppresses the equation number on a line of a derivation that
% you never need to refer back to.
\begin{align}
    \frac{T - T_\mathrm{env}}{T_0 - T_\mathrm{env}}
        &= \mathrm{e}^{-t/\tau} \notag \\
    t &= -\tau \ln \left( \frac{T - T_\mathrm{env}}{T_0 - T_\mathrm{env}} \right) .
    \label{eq:waiting}
\end{align}

Equations~\eqref{eq:cooling} and~\eqref{eq:waiting} are all the theory this report needs:
the first says what the curve looks like, the second says when to pick the cup up.

What counts as pleasant is of course a matter of taste.
This report uses

\[
    \operatorname{verdict}(T) =
    \begin{cases}
        \text{too hot}   & \text{if} \quad T > 65\,^{\circ}\mathrm{C} \\
        \text{drinkable} & \text{if} \quad 50\,^{\circ}\mathrm{C} \leq T \leq 65\,^{\circ}\mathrm{C} \\
        \text{too cold}  & \text{if} \quad T < 50\,^{\circ}\mathrm{C}
    \end{cases}
\]

% =====================================================================
\section{Setup}\label{sec:setup}

Each vessel starts with 300 ml of water at 85 degrees Celsius.
A thermometer sits in the centre of the liquid, clear of the walls.
The temperature is taken every 30 seconds for one hour.

The room stays closed and unventilated throughout, at 21 degrees Celsius.

The procedure is as follows:

\begin{enumerate}
    \item Heat 300 ml of water to 85 degrees Celsius.
    \item Pour it into the vessel and start the clock.
    \item Suspend the thermometer in the centre of the liquid.
    \item Record the temperature every 30 seconds for one hour.
\end{enumerate}

% TODO 4.1  Type the vessel list as a bare tabular, right here, so you
%           can see what a tabular on its own does to the page:
%
%             Label | Vessel     | Material       | Lid | tau
%             mug   | Mug        | glazed ceramic | no  | 24 min
%             glass | Glass      | borosilicate   | no  | 18 min
%             flask | Travel mug | steel, vacuum  | yes | 310 min
%
%           Five columns: four left, one right. Remember that &
%           separates columns and \\ ends a row -- this is the one
%           place \\ genuinely belongs. Put the labels in \texttt{}.

% TODO 4.2  Now wrap that tabular in a table float so it stops sitting
%           in the middle of your sentence:
%
%             \begin{table}[htbp]
%                 \centering
%                 \caption{The three vessels and their time constants.}
%                 \label{tab:vessels}
%                 ... your tabular ...
%             \end{table}
%
%           \caption goes BEFORE \label, always. Then replace your
%           \hline rules with \toprule, \midrule and \bottomrule from
%           booktabs, and compare the two PDFs side by side.
%           Load \usepackage{caption} too: without it, the gap between
%           the caption and the table is far too small.

% =====================================================================
\section{Results}\label{sec:results}

% TODO 4.3  Place the curves as a figure float. The file is
%           figures/cooling-curves.pdf -- note that you leave the
%           extension off:
%
%             \begin{figure}[htbp]
%                 \centering
%                 \includegraphics[width=0.8\textwidth]{figures/cooling-curves}
%                 \caption{Cooling curves of the three vessels.}
%                 \label{fig:cooling}
%             \end{figure}

% TODO 4.4  Write a short paragraph of results underneath, and refer to
%           both the figure and the table by number:
%
%             Figure~\ref{fig:cooling} ...   Table~\ref{tab:vessels} ...
%
%           Mind the ~ : it is a space that LaTeX is not allowed to
%           break, so "Figure" and its number never end up on different
%           lines. Mention that the mug is drinkable from about 9 to 19
%           minutes, and that the flask is still too hot nearly two
%           hours later.

% =====================================================================
\section{Conclusion}\label{sec:conclusion}

% TODO 4.5  Write two sentences of conclusion. Cite one of the sources
%           sitting in references.bib -- but misspell the key on
%           purpose the first time:
%
%             \cite{halliday2019}
%
%           Add \printbibliography immediately before \end{document},
%           then compile and look at two things: what appears in the
%           PDF where the citation should be, and what the log says
%           about it. THEN fix the key to \cite{halliday2018} and
%           compile again.
%           This is the error you will hit most often in a real report,
%           and now you know what it looks like.

% TODO 4.6  Open references.bib and follow the TODO at the bottom of
%           that file: add a third source of your own and cite it.

\end{document}

% =====================================================================
%  BLOCK 4 -- TABLES, FIGURES AND CITATIONS
% =====================================================================
%
%  The thread running through this block: you never type a number that
%  LaTeX can work out for you. Not a section number, not an equation
%  number, not a figure number, not a reference number. You give the
%  thing a \label and you ask for it by name.
%
%  AND THE PACKAGES WORTH ADDING NEXT
%  Including the tidier way of writing numbers with units that was
%  promised back in block 3:
%
%    \usepackage{siunitx}    then  \qty{65}{\degreeCelsius}
%                                  \qty{300}{\milli\litre}
%                                  \qty{24}{\minute}
%            It gets the thin space, the font and the minus signs
%            right, every time, and you can change the convention for
%            the whole document in one line.
%
%    \usepackage{microtype}  quietly better justification, no effort
%
%    \usepackage{hyperref}   every \ref, \eqref, \cite and table-of-
%                            contents entry becomes clickable.
%                            Load this one LAST of all your packages.
%
%  IF YOU FINISH EARLY
%    * Add siunitx and convert the numbers in your table and your cases
%      block to \qty{}{}.
%    * Split the document up: make a sections/ directory, move the
%      Theory section into sections/theory.tex, and pull it back in
%      with \input{sections/theory}. Worth doing once your document is
%      longer than this one.
% =====================================================================
